\documentclass[11pt,letterpaper]{article}

\usepackage[margin=2.2cm]{geometry}
\usepackage{amsmath}
\usepackage{amssymb}
\usepackage{array}
\usepackage{booktabs}

\title{Memoria de calculo detallada\\Proyecto IE0413}
\author{Daniel Gerardo Madriz Chavarr\'ia C24478\\
Fabricio Torres Jim\'enez C27803\\
Kenneth Daniel Vizca\'ino Jim\'enez C09331}
\date{I ciclo 2026}

\begin{document}

\maketitle

\section*{Nota de uso}

La realimentacion recibida indica que los calculos pueden presentarse en
\LaTeX, siempre que se muestre cada paso. Por eso se escriben las ecuaciones,
los despejes y las sustituciones numericas. La onda cuadrada se genera con un
\texttt{NE555}; las etapas con amplificador operacional se simulan en LTspice
con el macromodelo no ideal \texttt{LF353/NS}.

\section{Datos del grupo}

El enunciado pide ordenar alfabeticamente los primeros apellidos:

\[
\text{Madriz}, \qquad \text{Torres}, \qquad \text{Vizca\'ino}.
\]

Por tanto, los integrantes quedan en este orden:

\[
\begin{array}{c|c|c}
\text{Orden} & \text{Integrante} & \text{Carne} \\
\hline
1 & \text{Daniel Gerardo Madriz Chavarr\'ia} & \text{C24478} \\
2 & \text{Fabricio Torres Jim\'enez} & \text{C27803} \\
3 & \text{Kenneth Daniel Vizca\'ino Jim\'enez} & \text{C09331}
\end{array}
\]

El primer digito no nulo del primer carne es:

\[
\text{C24478}\quad \Rightarrow \quad X=2.
\]

El primer digito no nulo del segundo carne es:

\[
\text{C27803}\quad \Rightarrow \quad Y=2.
\]

El primer digito no nulo del tercer carne es:

\[
\text{C09331}\quad \Rightarrow \quad Z=9.
\]

La frecuencia del oscilador central es:

\[
f_0=X,YZ~\text{kHz}.
\]

En forma decimal:

\[
f_0=
\left(
X+\frac{Y}{10}+\frac{Z}{100}
\right)1000~\text{Hz}.
\]

Sustituyendo:

\[
f_0=
\left(
2+\frac{2}{10}+\frac{9}{100}
\right)1000.
\]

\[
f_0=
\left(
2+0.2+0.09
\right)1000.
\]

\[
f_0=2.29(1000)=2290~\text{Hz}.
\]

Por tanto:

\[
\boxed{f_0=2.29~\text{kHz}}.
\]

\section{Oscilador central de onda cuadrada}

\subsection{Conexion del NE555}

Se usa un \texttt{NE555} en modo astable tipo Schmitt. Las entradas
\texttt{TRIG} y \texttt{THRES} se conectan al mismo capacitor. La salida
\texttt{OUT} carga y descarga ese capacitor por medio de una sola resistencia
\(R_{osc}\). La entrada \texttt{RESET} se conecta a \(V_{CC}\) y la terminal
\texttt{CONT} se desacopla con un capacitor pequeno.

El \texttt{NE555} cambia de estado cuando el capacitor cruza los niveles
internos:

\[
V_L=\frac{1}{3}V_{CC},
\qquad
V_H=\frac{2}{3}V_{CC}.
\]

Como el mismo resistor se usa para la carga y la descarga, se obtiene un ciclo
de trabajo cercano al \(50\%\).

\subsection{Tiempo de carga}

Cuando la salida esta en alto, el capacitor se carga hacia \(V_{CC}\). La
ecuacion de carga es:

\[
v_C(t)=V_{CC}+\left[v_C(0)-V_{CC}\right]e^{-t/(RC)}.
\]

Al inicio de la carga:

\[
v_C(0)=\frac{1}{3}V_{CC}.
\]

Entonces:

\[
v_C(t)=V_{CC}+
\left(\frac{1}{3}V_{CC}-V_{CC}\right)e^{-t/(RC)}.
\]

\[
v_C(t)=V_{CC}-\frac{2}{3}V_{CC}e^{-t/(RC)}.
\]

El cambio ocurre cuando:

\[
v_C(t_H)=\frac{2}{3}V_{CC}.
\]

Sustituyendo:

\[
\frac{2}{3}V_{CC}
=
V_{CC}-\frac{2}{3}V_{CC}e^{-t_H/(RC)}.
\]

Se divide entre \(V_{CC}\):

\[
\frac{2}{3}
=
1-\frac{2}{3}e^{-t_H/(RC)}.
\]

Se despeja:

\[
\frac{2}{3}e^{-t_H/(RC)}
=
1-\frac{2}{3}
=
\frac{1}{3}.
\]

\[
e^{-t_H/(RC)}
=
\frac{1/3}{2/3}
=
\frac{1}{2}.
\]

Aplicando logaritmo natural:

\[
-\frac{t_H}{RC}=\ln\left(\frac{1}{2}\right).
\]

\[
t_H=RC\ln(2).
\]

\subsection{Tiempo de descarga}

Cuando la salida esta en bajo, el capacitor se descarga hacia \(0~\text{V}\).
La ecuacion es:

\[
v_C(t)=v_C(0)e^{-t/(RC)}.
\]

Al inicio de la descarga:

\[
v_C(0)=\frac{2}{3}V_{CC}.
\]

Por tanto:

\[
v_C(t)=\frac{2}{3}V_{CC}e^{-t/(RC)}.
\]

El cambio ocurre cuando:

\[
v_C(t_L)=\frac{1}{3}V_{CC}.
\]

Sustituyendo:

\[
\frac{1}{3}V_{CC}
=
\frac{2}{3}V_{CC}e^{-t_L/(RC)}.
\]

Se divide entre \(V_{CC}\):

\[
\frac{1}{3}
=
\frac{2}{3}e^{-t_L/(RC)}.
\]

\[
e^{-t_L/(RC)}
=
\frac{1}{2}.
\]

Entonces:

\[
t_L=RC\ln(2).
\]

Como:

\[
t_H=t_L,
\]

el ciclo de trabajo ideal es:

\[
D=\frac{t_H}{t_H+t_L}
=
\frac{1}{2}
=50\%.
\]

El periodo total es:

\[
T=t_H+t_L.
\]

\[
T=2RC\ln(2).
\]

Como:

\[
f_0=\frac{1}{T},
\]

entonces:

\[
f_0=\frac{1}{2RC\ln(2)}.
\]

Se despeja:

\[
R=\frac{1}{2f_0C\ln(2)}.
\]

\subsection{Sustitucion numerica}

Se escoge:

\[
C_{osc}=10~\text{nF}=10\times 10^{-9}~\text{F}.
\]

Entonces:

\[
R_{osc,teo}=
\frac{1}
{2(2290)(10\times 10^{-9})\ln(2)}.
\]

Primero:

\[
2(2290)(10\times 10^{-9})
=4.58\times 10^{-5}.
\]

Luego:

\[
(4.58\times 10^{-5})\ln(2)
=
(4.58\times 10^{-5})(0.6931).
\]

\[
(4.58\times 10^{-5})(0.6931)
=3.1746\times 10^{-5}.
\]

Finalmente:

\[
R_{osc,teo}
=
\frac{1}{3.1746\times 10^{-5}}
=31499.9~\Omega.
\]

\[
R_{osc,teo}\approx 31.5~\text{k}\Omega.
\]

En la simulacion se usa:

\[
K_{trim,555}=0.99852.
\]

Entonces:

\[
R_{osc,LTspice}=K_{trim,555}R_{osc,teo}.
\]

\[
R_{osc,LTspice}=(0.99852)(31499.9~\Omega).
\]

\[
R_{osc,LTspice}=31453.3~\Omega.
\]

\[
\boxed{R_{osc,LTspice}\approx 31.46~\text{k}\Omega}.
\]

\section{Onda diente de sierra positiva}

El enunciado pide que la diente de sierra se repita tanto en el semiciclo
positivo como en el negativo de la onda cuadrada. Por eso:

\[
f_{saw}=2f_0.
\]

Sustituyendo:

\[
f_{saw}=2(2290)=4580~\text{Hz}.
\]

\[
\boxed{f_{saw}=4.58~\text{kHz}}.
\]

El periodo correspondiente es:

\[
T_{saw}=\frac{1}{f_{saw}}.
\]

\[
T_{saw}=\frac{1}{4580}.
\]

\[
T_{saw}=2.1834\times 10^{-4}~\text{s}.
\]

\[
T_{saw}=0.2183~\text{ms}.
\]

La etapa propuesta carga un capacitor desde \(+15~\text{V}\) por medio de una
resistencia \(R_{charge}\). En cada flanco de la onda cuadrada se activa un
transistor de descarga para regresar rapidamente el capacitor a \(0~\text{V}\).

La ecuacion de carga del capacitor es:

\[
v_C(t)=V_{CC}\left(1-e^{-t/(RC)}\right).
\]

Se desea que al final de cada periodo de la diente de sierra:

\[
v_C(T_{saw})=V_{saw,pk}.
\]

Sustituyendo:

\[
V_{saw,pk}
=
V_{CC}\left(1-e^{-T_{saw}/(RC)}\right).
\]

Se divide entre \(V_{CC}\):

\[
\frac{V_{saw,pk}}{V_{CC}}
=
1-e^{-T_{saw}/(RC)}.
\]

Se despeja el exponencial:

\[
e^{-T_{saw}/(RC)}
=
1-\frac{V_{saw,pk}}{V_{CC}}.
\]

Aplicando logaritmo natural:

\[
-\frac{T_{saw}}{RC}
=
\ln\left(1-\frac{V_{saw,pk}}{V_{CC}}\right).
\]

Se despeja \(R\):

\[
R=
\frac{-T_{saw}}
{C\ln\left(1-\frac{V_{saw,pk}}{V_{CC}}\right)}.
\]

Se escoge:

\[
V_{CC}=15~\text{V},
\qquad
V_{saw,pk}=5~\text{V},
\qquad
C_{saw}=10~\text{nF}.
\]

Entonces:

\[
R_{charge}
=
\frac{-2.1834\times 10^{-4}}
{(10\times 10^{-9})\ln\left(1-\frac{5}{15}\right)}.
\]

Primero:

\[
1-\frac{5}{15}
=
1-\frac{1}{3}
=
\frac{2}{3}.
\]

\[
\ln\left(\frac{2}{3}\right)=-0.4055.
\]

Luego:

\[
(10\times 10^{-9})(-0.4055)
=
-4.055\times 10^{-9}.
\]

Finalmente:

\[
R_{charge}
=
\frac{-2.1834\times 10^{-4}}
{-4.055\times 10^{-9}}.
\]

\[
R_{charge}=53849~\Omega.
\]

\[
\boxed{R_{charge}\approx 53.8~\text{k}\Omega}.
\]

En la simulacion se usa:

\[
\boxed{R_{charge}=53.6~\text{k}\Omega},
\]

valor practico cercano al calculado. La onda sigue siendo exclusivamente
positiva porque el transistor de descarga lleva el capacitor hacia tierra, no
hacia tension negativa.

Si se aproxima la rampa por una pendiente promedio, la corriente media
equivalente es:

\[
I_{saw,eq}=C_{saw}\frac{V_{saw,pk}}{T_{saw}}.
\]

\[
I_{saw,eq}=
(10\times 10^{-9})\frac{5}{2.1834\times 10^{-4}}.
\]

\[
I_{saw,eq}=229~\mu\text{A}.
\]

\section{Onda triangular}

La frecuencia de salida triangular debe ser:

\[
f_{tri}=\frac{f_0}{2}.
\]

Sustituyendo:

\[
f_{tri}=\frac{2290}{2}=1145~\text{Hz}.
\]

\[
\boxed{f_{tri}=1.145~\text{kHz}}.
\]

El periodo es:

\[
T_{tri}=\frac{1}{f_{tri}}=\frac{1}{1145}.
\]

\[
T_{tri}=8.7336\times 10^{-4}~\text{s}.
\]

\[
T_{tri}=0.8734~\text{ms}.
\]

El divisor por dos entrega una senal entre \(0~\text{V}\) y \(2~\text{V}\). El
cambiador de nivel resta \(1~\text{V}\):

\[
v_{drv}=v_{div}-1~\text{V}.
\]

Si \(v_{div}=0~\text{V}\):

\[
v_{drv}=0-1=-1~\text{V}.
\]

Si \(v_{div}=2~\text{V}\):

\[
v_{drv}=2-1=+1~\text{V}.
\]

Entonces:

\[
V_{drv}=1~\text{V}.
\]

Para el integrador inversor:

\[
\left|\frac{dv_o}{dt}\right|=
\frac{V_{drv}}{R_{tri}C_{tri}}.
\]

La triangular cambia de \(-V_{tri,pk}\) a \(+V_{tri,pk}\), por lo que:

\[
\Delta V=2V_{tri,pk}.
\]

Ese cambio ocurre durante medio periodo triangular:

\[
\frac{T_{tri}}{2}
=
\frac{1}{2f_{tri}}
=
\frac{1}{2(f_0/2)}
=
\frac{1}{f_0}.
\]

Entonces la pendiente requerida es:

\[
\left|\frac{dv_o}{dt}\right|
=
\frac{2V_{tri,pk}}{1/f_0}
=
2V_{tri,pk}f_0.
\]

Se igualan ambas pendientes:

\[
\frac{V_{drv}}{R_{tri}C_{tri}}
=
2V_{tri,pk}f_0.
\]

Se despeja:

\[
R_{tri}=
\frac{V_{drv}}
{2V_{tri,pk}f_0C_{tri}}.
\]

Se escoge:

\[
V_{tri,pk}=2~\text{V},
\qquad
C_{tri}=10~\text{nF}.
\]

Sustituyendo:

\[
R_{tri}=
\frac{1}
{2(2)(2290)(10\times 10^{-9})}.
\]

\[
2(2)(2290)=9160.
\]

\[
9160(10\times 10^{-9})=9.16\times 10^{-5}.
\]

\[
R_{tri}=\frac{1}{9.16\times 10^{-5}}
=10917.0~\Omega.
\]

\[
\boxed{R_{tri}\approx 10.9~\text{k}\Omega}.
\]

\section{Seleccion de la onda para obtener la senoidal}

La onda cuadrada simetrica contiene armonicos impares:

\[
v_{cuad}(t)=
\frac{4A}{\pi}
\left[
\sin(\omega_0 t)
+\frac{1}{3}\sin(3\omega_0 t)
+\frac{1}{5}\sin(5\omega_0 t)
+\cdots
\right].
\]

Sus armonicos decrecen como:

\[
\frac{1}{n}.
\]

La onda diente de sierra contiene armonicos pares e impares:

\[
v_{saw}(t)=
\frac{2A}{\pi}
\left[
\sin(\omega_0 t)
-\frac{1}{2}\sin(2\omega_0 t)
+\frac{1}{3}\sin(3\omega_0 t)
-\cdots
\right].
\]

Sus armonicos tambien decrecen como:

\[
\frac{1}{n}.
\]

La onda triangular simetrica contiene solamente armonicos impares:

\[
v_{tri}(t)=
\frac{8A}{\pi^2}
\left[
\sin(\omega_0 t)
-\frac{1}{3^2}\sin(3\omega_0 t)
+\frac{1}{5^2}\sin(5\omega_0 t)
-\cdots
\right].
\]

Sus armonicos decrecen como:

\[
\frac{1}{n^2}.
\]

Comparando el tercer armonico:

\[
\text{cuadrada:}\quad \frac{1}{3}=0.333.
\]

\[
\text{triangular:}\quad \frac{1}{3^2}=\frac{1}{9}=0.111.
\]

La relacion es:

\[
\frac{1/9}{1/3}=\frac{1}{3}.
\]

El tercer armonico de la triangular es tres veces menor que el de la cuadrada
antes de filtrar. Ademas, la triangular no tiene armonicos pares. Por eso se
escoge la onda triangular como entrada del filtro pasabajas.

\section{Filtro pasabajas de orden 5}

\subsection{Frecuencia de corte}

La fundamental de la triangular es:

\[
f_{tri}=1145~\text{Hz}.
\]

Se escoge:

\[
f_c=1.25f_{tri}.
\]

\[
f_c=1.25(1145)=1431.25~\text{Hz}.
\]

\[
\boxed{f_c=1.43125~\text{kHz}}.
\]

El primer armonico no deseado de la triangular es el tercero:

\[
f_3=3f_{tri}.
\]

\[
f_3=3(1145)=3435~\text{Hz}.
\]

La relacion respecto a \(f_c\) es:

\[
\frac{f_3}{f_c}
=
\frac{3435}{1431.25}
=2.4.
\]

Por tanto, el tercer armonico queda en \(2.4f_c\).

\subsection{Tipo de respuesta}

El tipo de respuesta depende de \(Z\):

\[
\begin{array}{c|c}
Z & \text{Respuesta} \\
\hline
1,2,3 & \text{Bessel}\\
4,5,6 & \text{Butterworth}\\
7,8,9 & \text{Chebyshev}
\end{array}
\]

Como:

\[
Z=9,
\]

se usa:

\[
\boxed{\text{Chebyshev tipo I de 0.5 dB}}.
\]

\subsection{Estructura del filtro}

El orden total es:

\[
5=1+2+2.
\]

Se implementa con:

\begin{itemize}
  \item una etapa de primer orden,
  \item una etapa Sallen-Key de segundo orden,
  \item una segunda etapa Sallen-Key de segundo orden.
\end{itemize}

Para la etapa de primer orden:

\[
f_p=\frac{1}{2\pi RC}.
\]

Se despeja:

\[
R=\frac{1}{2\pi f_p C}.
\]

Para cada Sallen-Key se usa:

\[
R_1=R_2=R,
\qquad
C_1=C_2=C.
\]

Entonces:

\[
f_n=\frac{1}{2\pi RC}.
\]

Se despeja:

\[
R=\frac{1}{2\pi f_n C}.
\]

La ganancia no inversora de la etapa es:

\[
K=1+\frac{R_f}{R_g}.
\]

Con componentes iguales en Sallen-Key:

\[
Q=\frac{1}{3-K}.
\]

Se despeja \(K\):

\[
Q(3-K)=1.
\]

\[
3Q-KQ=1.
\]

\[
KQ=3Q-1.
\]

\[
K=\frac{3Q-1}{Q}.
\]

\[
K=3-\frac{1}{Q}.
\]

\subsection{Coeficientes Chebyshev}

Para Chebyshev de 0.5 dB se usan:

\[
\begin{array}{c|c}
\text{Parametro} & \text{Valor} \\
\hline
f_{p1}/f_c & 0.3623\\
f_{n1}/f_c & 0.6905\\
Q_1 & 1.1778\\
f_{n2}/f_c & 1.0177\\
Q_2 & 4.5450
\end{array}
\]

\subsection{Etapa de primer orden}

\[
f_{p1}=0.3623f_c.
\]

\[
f_{p1}=0.3623(1431.25).
\]

\[
f_{p1}=518.57~\text{Hz}.
\]

Se escoge:

\[
C_f=10~\text{nF}.
\]

Entonces:

\[
R_{p1}=
\frac{1}{2\pi f_{p1}C_f}.
\]

\[
R_{p1}=
\frac{1}
{2\pi(518.57)(10\times 10^{-9})}.
\]

\[
R_{p1}=30691.1~\Omega.
\]

\[
\boxed{R_{p1}\approx 30.7~\text{k}\Omega}.
\]

\subsection{Primera etapa Sallen-Key}

Para la primera seccion Chebyshev:

\[
Q_1=1.1778.
\]

La ganancia requerida es:

\[
K_1=3-\frac{1}{Q_1}.
\]

\[
K_1=3-\frac{1}{1.1778}.
\]

\[
\frac{1}{1.1778}=0.8490.
\]

\[
K_1=3-0.8490=2.1510.
\]

La frecuencia natural es:

\[
f_{n1}=0.6905f_c.
\]

\[
f_{n1}=0.6905(1431.25).
\]

\[
f_{n1}=988.25~\text{Hz}.
\]

La resistencia es:

\[
R_{SK1}=
\frac{1}{2\pi f_{n1}C_f}.
\]

\[
R_{SK1}=
\frac{1}
{2\pi(988.25)(10\times 10^{-9})}.
\]

\[
R_{SK1}=16104.7~\Omega.
\]

\[
\boxed{R_{SK1}\approx 16.1~\text{k}\Omega}.
\]

Para implementar la ganancia:

\[
K_1=1+\frac{R_{f1}}{R_{g1}}.
\]

Se escoge:

\[
R_{g1}=10~\text{k}\Omega.
\]

Se despeja:

\[
R_{f1}=(K_1-1)R_{g1}.
\]

\[
R_{f1}=(2.1510-1)(10~\text{k}\Omega).
\]

\[
R_{f1}=11.51~\text{k}\Omega.
\]

\subsection{Segunda etapa Sallen-Key}

Para la segunda seccion Chebyshev:

\[
Q_2=4.5450.
\]

La ganancia requerida es:

\[
K_2=3-\frac{1}{Q_2}.
\]

\[
K_2=3-\frac{1}{4.5450}.
\]

\[
\frac{1}{4.5450}=0.2200.
\]

\[
K_2=3-0.2200=2.7800.
\]

La frecuencia natural es:

\[
f_{n2}=1.0177f_c.
\]

\[
f_{n2}=1.0177(1431.25).
\]

\[
f_{n2}=1456.63~\text{Hz}.
\]

La resistencia es:

\[
R_{SK2}=
\frac{1}{2\pi f_{n2}C_f}.
\]

\[
R_{SK2}=
\frac{1}
{2\pi(1456.63)(10\times 10^{-9})}.
\]

\[
R_{SK2}=10926.2~\Omega.
\]

\[
\boxed{R_{SK2}\approx 10.9~\text{k}\Omega}.
\]

Para la ganancia:

\[
K_2=1+\frac{R_{f2}}{R_{g2}}.
\]

Se escoge:

\[
R_{g2}=10~\text{k}\Omega.
\]

Entonces:

\[
R_{f2}=(K_2-1)R_{g2}.
\]

\[
R_{f2}=(2.7800-1)(10~\text{k}\Omega).
\]

\[
R_{f2}=17.80~\text{k}\Omega.
\]

\subsection{Atenuador de entrada}

La ganancia total de las dos etapas Sallen-Key es:

\[
K_T=K_1K_2.
\]

\[
K_T=(2.1510)(2.7800).
\]

\[
K_T=5.9796.
\]

Para no amplificar excesivamente la triangular, se coloca un divisor de entrada
con ganancia:

\[
A=\frac{1}{K_T}.
\]

El divisor cumple:

\[
A=\frac{R_{bot}}{R_{top}+R_{bot}}.
\]

Se despeja \(R_{top}\):

\[
A(R_{top}+R_{bot})=R_{bot}.
\]

\[
AR_{top}+AR_{bot}=R_{bot}.
\]

\[
AR_{top}=R_{bot}(1-A).
\]

\[
R_{top}=\frac{R_{bot}(1-A)}{A}.
\]

Como \(A=1/K_T\):

\[
R_{top}=R_{bot}(K_T-1).
\]

Se escoge:

\[
R_{bot}=100~\text{k}\Omega.
\]

Entonces:

\[
R_{top}=(100~\text{k}\Omega)(5.9796-1).
\]

\[
R_{top}=497.96~\text{k}\Omega.
\]

\[
\boxed{R_{top}\approx 498~\text{k}\Omega}.
\]

\subsection{Verificacion del rechazo del tercer armonico}

En la simulacion AC del filtro se obtiene:

\[
|H(f_{tri})|=-11.12~\text{dB}.
\]

\[
|H(3f_{tri})|=-62.29~\text{dB}.
\]

Por tanto, el filtro atenua el tercer armonico, respecto a la fundamental, en:

\[
62.29-11.12=51.17~\text{dB}.
\]

Ademas, para la onda triangular:

\[
\frac{A_3}{A_1}=\frac{1}{3^2}=\frac{1}{9}.
\]

En decibeles:

\[
20\log_{10}\left(\frac{1}{9}\right)=-19.08~\text{dB}.
\]

Entonces, a la salida del filtro, el tercer armonico queda aproximadamente:

\[
51.17+19.08=70.25~\text{dB}
\]

por debajo de la fundamental. Esto justifica la seleccion de la triangular y
del filtro pasabajas de orden 5.

\section{Resumen numerico}

\[
\begin{array}{c|c}
\text{Parametro} & \text{Valor} \\
\hline
X,Y,Z & 2,2,9\\
f_0 & 2290~\text{Hz}\\
f_{saw} & 4580~\text{Hz}\\
f_{tri} & 1145~\text{Hz}\\
f_c & 1431.25~\text{Hz}\\
R_{osc,LTspice} & 31.46~\text{k}\Omega\\
R_{charge} & 53.6~\text{k}\Omega\\
R_{tri} & 10.9~\text{k}\Omega\\
R_{p1} & 30.7~\text{k}\Omega\\
R_{SK1} & 16.1~\text{k}\Omega\\
R_{f1} & 11.5~\text{k}\Omega\\
R_{SK2} & 10.9~\text{k}\Omega\\
R_{f2} & 17.8~\text{k}\Omega\\
R_{top} & 498~\text{k}\Omega\\
R_{bot} & 100~\text{k}\Omega
\end{array}
\]

\end{document}
